%
% Capítulo 1
%
\chapter{Introdução} \label{cap:intro}

A organização de horários entre professores e alunos é uma tarefa que, em muitos contextos, continua a ser realizada
de forma manual, recorrendo a trocas sucessivas de mensagens, folhas e imagens ou registos informais.
Este processo tende a ser demorado, pouco flexível e propenso a erros, especialmente quando existem múltiplos alunos,
disponibilidades distintas e restrições específicas que precisam de ser respeitadas. A dificuldade aumenta quando se
pretende encontrar uma distribuição equilibrada de sessões, compatível com a disponibilidade do professor e com as
preferências ou limitações dos alunos.

A motivação para este projeto surgiu de uma experiência concreta das autoras no contexto de um centro de estudos, onde
ambas desempenham funções de explicadoras. Neste contexto, a gestão dos horários dos alunos era realizada manualmente e
exigia reajustes frequentes sempre que entrava um novo aluno ou surgiam alterações de disponibilidade. Na prática, isso
implicava refazer sucessivamente os horários já definidos, tentar encontrar novos encaixes e, muitas vezes, lidar com a
dificuldade de conciliar todos os alunos de forma equilibrada e sustentável. Esta experiência permitiu identificar de
forma clara a necessidade de uma solução que apoiasse este processo de forma mais eficiente e estruturada.

Neste contexto, surgiu a proposta do projeto \textit{Sistema de Planeamento de Horários}, cujo objetivo consiste no
desenvolvimento de uma solução capaz de apoiar a gestão de horários entre professores e alunos.
A solução permite recolher disponibilidades, definir restrições de agendamento e construir propostas de horário de forma
estruturada, reduzindo o esforço manual associado a este processo e tornando a gestão de sessões mais clara e eficiente.

%
% Secção 1.1
%
\section{Motivação e Objetivos} \label{sec11}

A motivação principal para o desenvolvimento deste projeto prende-se com a necessidade de criar uma ferramenta capaz de simplificar a coordenação de horários em cenários onde um professor trabalha com vários alunos e necessita de conciliar diferentes disponibilidades ao longo da semana. Em vez de depender exclusivamente de planeamento manual, a solução proposta procura automatizar parte desse trabalho, permitindo que os dados relevantes sejam recolhidos de forma consistente e utilizados
para construir propostas de horário com base em regras bem definidas.

O objetivo central do projeto consiste, assim, em disponibilizar um sistema que permita:

\begin{itemize}
    \item registar professores e alunos;
    \item associar alunos a um professor;
    \item recolher as disponibilidades semanais de professores e alunos;
    \item definir restrições relevantes para o agendamento;
    \item construir propostas de horário;
    \item persistir aulas concretas com datas reais;
    \item apoiar a gestão operacional das aulas, incluindo consulta, remarcação, cancelamento e marcação de presenças;
    \item apresentar o horário resultante de forma legível, organizada e acessível.
\end{itemize}

Numa perspetiva mais ampla, o projeto pretende também servir como base para a evolução de funcionalidades futuras, nomeadamente, ao nível da gestão de salas e do suporte à interpretação automática de disponibilidades a partir de texto, imagem ou áudio.

%
% Secção 1.2
%
\section{Visão Geral da Solução} \label{sec12}

A solução desenvolvida assenta numa arquitetura composta por uma aplicação Android, responsável pela interação com o utilizador e por um backend com API REST, responsável pela persistência central de dados, pela execução da lógica principal do sistema e pela exposição dos serviços necessários ao frontend. No frontend foi utilizada a linguagem Kotlin com Jetpack Compose para a construção da interface, enquanto no backend foi adotado Ktor para a implementação da API, em conjunto com PostgreSQL para persistência remota.

A aplicação distingue dois perfis principais de utilizador: professor e aluno. Cada um destes perfis dispõe de um conjunto próprio de funcionalidades e ecrãs adaptados ao seu papel no sistema. O professor pode definir a sua disponibilidade semanal, configurar restrições de agendamento, consultar os alunos associados, criar propostas de horário, confirmar essas propostas, persistir aulas concretas com datas reais e gerir o ciclo de vida das aulas já criadas. O aluno pode escolher o professor com quem pretende trabalhar, definir a sua disponibilidade, indicar a sua carga horária semanal pretendida e consultar as aulas que lhe estão atribuídas.

A partir da informação recolhida, o sistema constrói propostas de horário semanal compatíveis com as disponibilidades registadas e com as restrições configuradas. Após a aceitação da proposta pelo professor, essas sessões são convertidas em aulas concretas com datas reais, recorrência semanal e persistência na base de dados. Sobre essas aulas, o sistema suporta ainda operações como consulta de histórico, marcação de presenças, cancelamento, remarcação e envio de notificações por email aos alunos.

Como complemento prático, o sistema integra também o Google Calendar do professor, permitindo registar automaticamente no calendário os eventos correspondentes às aulas aceites. Desta forma, a solução aproxima-se de um cenário real de utilização, não se limitando à geração abstrata de horários, mas apoiando também a sua gestão operacional ao longo do tempo.

%
% Secção 1.3
%
\section{Organização do Documento} \label{sec13}

O presente relatório encontra-se organizado da seguinte forma. No Capítulo~\ref{cap:enquadramento} é apresentado o enquadramento do problema, a motivação do projeto e o contexto em que a solução se insere.
No Capítulo~\ref{cap:arquitetura} é descrita a organização global do sistema, incluindo os principais componentes e a sua interação.
Seguidamente, são apresentados os detalhes de implementação mais relevantes, tanto ao nível do frontend como do backend, bem como os principais aspetos da API, da persistência de dados e da lógica de criação de horários.
Posteriormente, é descrita a estratégia de testes e validação aplicada ao sistema.
É também apresentado um estudo exploratório sobre a utilização de modelos multimodais para interpretação de disponibilidades a partir de diferentes tipos de input.
Por fim, é feito um balanço do estado final do sistema, identificando as funcionalidades concluídas, as limitações atuais, os principais resultados obtidos e as conclusões fundamentais do trabalho realizado.

